% ══════════════════════════════════════════════════════════════════════════════
%                     کتاب: راست یا چپ؛ میانه کجاست؟
%                          نویسنده: مهدی سالم
%                           شهر: ریچموندهیل
% ══════════════════════════════════════════════════════════════════════════════
%                            فایل اصلی: main.tex
%                    کامپایلر: XeLaTeX | فونت: Vazirmatn
% ══════════════════════════════════════════════════════════════════════════════

\documentclass[12pt, a4paper, twoside, openany]{book}

% ══════════════════════════════════════════════════════════════════════════════
%                              پکیج‌های پایه
% ══════════════════════════════════════════════════════════════════════════════

\usepackage[
    top=2.5cm,
    bottom=2.5cm,
    left=2cm,
    right=2.5cm,
    headheight=36pt,
    headsep=1cm,
    footskip=1.5cm,
    bindingoffset=0.5cm]{geometry}

% ریاضی و نمادها
\usepackage{amsmath, amssymb, amsfonts}

% گرافیک و رنگ
\usepackage[dvipsnames, svgnames, x11names, table]{xcolor}
\usepackage{graphicx}
\usepackage{tikz}
\usepackage{pgfplots}
\pgfplotsset{compat=1.18}

% کتابخانه‌های TikZ
\usetikzlibrary{
    shapes.geometric,
    shapes.symbols,
    shapes.arrows,
    shapes.callouts,
    shapes.misc,
    arrows.meta,
    positioning,
    calc,
    fit,
    backgrounds,
    decorations.pathreplacing,
    decorations.pathmorphing,
    decorations.markings,
    decorations.text,
    patterns,
    shadows,
    shadows.blur,
    fadings,
    chains,
    matrix,
    mindmap,
    calendar,
    quotes,
    angles,
    intersections,
    through,
    spy,
    turtle,
    lindenmayersystems,
    folding,
    patterns.meta,
    bending
}

% جداول پیشرفته
\usepackage{booktabs}
\usepackage{longtable}
\usepackage{multirow}
\usepackage{makecell}
\usepackage{colortbl}
\usepackage{array}
\usepackage{tabularx}

% کادرها و باکس‌ها
\usepackage[most, skins, breakable]{tcolorbox}

% صفحات عرضی
\usepackage{pdflscape}
\usepackage{rotating}

% هدر و فوتر
\usepackage{fancyhdr}

% لینک‌ها
\usepackage{hyperref}
\hypersetup{
    colorlinks=true,
    linkcolor=RoyalBlue4,
    citecolor=ForestGreen,
    urlcolor=Maroon,
    pdfencoding=unicode,
    unicode=true,
    pdflang={fa-IR},
    bookmarks=true,
    bookmarksopen=true,
    pdfauthor={مهدی سالم},
    pdftitle={راست یا چپ؛ میانه کجاست؟},
    pdfsubject={اندیشه سیاسی},
    pdfkeywords={چپ، راست، سیاست، ایدئولوژی}
}

% شماره‌گذاری و فهرست
\usepackage{tocbibind}
\usepackage{titletoc}
\usepackage{titlesec}

% پاورقی
\usepackage[perpage, symbol*]{footmisc}

% نقل‌قول
\usepackage{csquotes}

% سایر
\usepackage{enumitem}
\usepackage{lettrine}
\usepackage{lipsum}
\usepackage{setspace}
\usepackage{float}
\usepackage{wrapfig}
\usepackage{caption}
\usepackage{subcaption}
\usepackage{needspace}

% در پریمبل اضافه کنید (قبل از xepersian):
\usepackage{svg}
\svgsetup{
	inkscapelatex=false,   % متن SVG را Inkscape رندر کند نه LaTeX
	inkscapeformat=pdf,
	inkscapearea=drawing
}

% ══════════════════════════════════════════════════════════════════════════════
%                   🎨 پالت رنگی اختصاصی کتاب
% ══════════════════════════════════════════════════════════════════════════════

% رنگ‌های اصلی طیف سیاسی
\definecolor{LeftRed}{RGB}{180, 30, 30}
\definecolor{LeftRedLight}{RGB}{220, 80, 80}
\definecolor{LeftRedDark}{RGB}{130, 20, 20}

\definecolor{RightBlue}{RGB}{30, 60, 140}
\definecolor{RightBlueLight}{RGB}{70, 100, 180}
\definecolor{RightBlueDark}{RGB}{20, 40, 100}

\definecolor{CenterGreen}{RGB}{50, 130, 70}
\definecolor{CenterGreenLight}{RGB}{90, 170, 110}
\definecolor{CenterGreenDark}{RGB}{30, 90, 50}

% رنگ‌های تأکیدی
\definecolor{GoldAccent}{RGB}{180, 150, 40}
\definecolor{GoldLight}{RGB}{220, 190, 80}
\definecolor{GoldDark}{RGB}{140, 110, 20}

\definecolor{PurpleAccent}{RGB}{100, 50, 120}
\definecolor{PurpleLight}{RGB}{150, 100, 170}

\definecolor{OrangeAccent}{RGB}{200, 100, 30}
\definecolor{OrangeLight}{RGB}{240, 150, 80}

% رنگ‌های خنثی
\definecolor{DarkGray}{RGB}{50, 50, 55}
\definecolor{MediumGray}{RGB}{100, 100, 105}
\definecolor{LightGray}{RGB}{200, 200, 205}
\definecolor{VeryLightGray}{RGB}{240, 240, 242}
\definecolor{PageBg}{RGB}{252, 251, 248}

% رنگ‌های کادرها
\definecolor{DefBoxBg}{RGB}{232, 245, 250}
\definecolor{DefBoxFrame}{RGB}{70, 130, 180}
\definecolor{QuoteBoxBg}{RGB}{235, 250, 240}
\definecolor{QuoteBoxFrame}{RGB}{60, 140, 90}
\definecolor{HistoryBoxBg}{RGB}{255, 250, 235}
\definecolor{HistoryBoxFrame}{RGB}{180, 150, 50}
\definecolor{DebateBoxBg}{RGB}{255, 240, 240}
\definecolor{DebateBoxFrame}{RGB}{180, 50, 50}
\definecolor{CompareBoxBg}{RGB}{245, 240, 255}
\definecolor{CompareBoxFrame}{RGB}{100, 60, 130}
\definecolor{PersonBoxBg}{RGB}{248, 248, 252}
\definecolor{PersonBoxFrame}{RGB}{80, 80, 120}
\definecolor{WarningBoxBg}{RGB}{255, 245, 230}
\definecolor{WarningBoxFrame}{RGB}{200, 120, 40}

% ── رنگ‌های اضافه برای عمق بصری ──────────────────────────────────────────
\definecolor{ChapterBandColor}{RGB}{25, 50, 120}   % نوار فصل
\definecolor{ChapterNumColor}{RGB}{220, 185, 60}   % شماره فصل طلایی
\definecolor{OrnamentColor}{RGB}{160, 130, 40}     % رنگ زیورها
\definecolor{PartBg}{RGB}{20, 40, 90}              % پس‌زمینه بخش
\definecolor{ShadeLine}{RGB}{180, 160, 100}        % خط‌های سایه‌دار

% Color aliases
\colorlet{leftred}{LeftRed}
\colorlet{rightblue}{RightBlue}
\colorlet{centergreen}{CenterGreen}
\colorlet{goldaccent}{GoldAccent}
\colorlet{violet}{PurpleAccent}
\colorlet{darkgray}{DarkGray}
\colorlet{lightgray}{LightGray}
\colorlet{verylightgray}{VeryLightGray}

% ══════════════════════════════════════════════════════════════════════════════
%              🖼️  زیور و تزئینات تکرارشونده (TikZ helpers)
% ══════════════════════════════════════════════════════════════════════════════

% ─── خط تزئینی طلایی ─────────────────────────────────────────────────────────
%  استفاده: \ornamentalrule   (در هر جا که \titlerule قدیمی بود)
% در پریمبل، ornamentalrule را کاملاً بازنویسی کنید:
\newcommand{\ornamentalrule}{%
  \begin{center}
  \begin{tikzpicture}[baseline=0pt]
    % خط چپ
    \draw[OrnamentColor, line width=0.8pt]
      (0,0) -- (0.28\linewidth, 0);
    % نقطهٔ گرد چپ
    \fill[OrnamentColor]
      (0.28\linewidth, 0) circle (2pt);
    % خط مرکزی ضخیم
    \draw[GoldAccent, line width=1.8pt]
      (0.28\linewidth, 0) -- (0.72\linewidth, 0);
    % لوزی مرکزی
    \node[diamond, fill=GoldAccent,
          inner sep=2.2pt, outer sep=0pt]
      at (0.5\linewidth, 0) {};
    % نقطهٔ گرد راست
    \fill[OrnamentColor]
      (0.72\linewidth, 0) circle (2pt);
    % خط راست
    \draw[OrnamentColor, line width=0.8pt]
      (0.72\linewidth, 0) -- (\linewidth, 0);
  \end{tikzpicture}
  \end{center}%
}

% bookrule هم بر اساس ornamentalrule تعریف شده — خودکار درست می‌شود
\newcommand{\bookrule}{%
  \vspace{4pt}\ornamentalrule\vspace{4pt}%
}

% ─── جداکننده بند (section ornament) ────────────────────────────────────────
\newcommand{\sectionornament}{%
  \begin{center}
  \begin{tikzpicture}
    \foreach \x in {-1.2,-0.6,0,0.6,1.2}{
      \filldraw[OrnamentColor!\x*30+60] (\x,0) circle (2pt);
    }
    \draw[OrnamentColor, line width=0.5pt] (-1.5,0)--(1.5,0);
  \end{tikzpicture}
  \end{center}%
}

% ─── corner ornament (برای صفحه عنوان / بخش) ─────────────────────────────────
\newcommand{\cornerornament}[1][OrnamentColor]{%
  \begin{tikzpicture}[remember picture, overlay]
    % گوشه بالا-راست
    \draw[#1, line width=1pt]
      ([xshift=-1.5cm, yshift=-1.5cm]current page.north east)
      -| ([xshift=-1.5cm, yshift=-1.5cm]current page.north east
          -| current page.north east)
         ([xshift=-1.5cm, yshift=-1.5cm]current page.north east)
      |- ([xshift=-1.5cm, yshift=-1.5cm]current page.north east
          |- current page.north east);
    % گوشه پایین-چپ
    \draw[#1, line width=1pt]
      ([xshift=1.5cm, yshift=1.5cm]current page.south west)
      -| ([xshift=1.5cm, yshift=1.5cm]current page.south west
          -| current page.south west)
         ([xshift=1.5cm, yshift=1.5cm]current page.south west)
      |- ([xshift=1.5cm, yshift=1.5cm]current page.south west
          |- current page.south west);
  \end{tikzpicture}%
}

% ══════════════════════════════════════════════════════════════════════════════
%                   📦 تعریف کادرهای رنگی (tcolorbox)
% ══════════════════════════════════════════════════════════════════════════════

% ─── تنظیم جهانی tcolorbox ─────────────────────────────────────────────────
% تمام کادرها از این پایه ارث می‌برند؛ تغییر اینجا = تغییر همه جا
\tcbset{
  bookbase/.style={
    enhanced,
    breakable,
    arc=3mm,
    boxrule=0.8pt,
    fonttitle=\bfseries,
    attach boxed title to top right={yshift=-2.5mm, xshift=-5mm},
    boxed title style={
      rounded corners, % was =2mm
      boxrule=0pt,
      interior style={
        left color=#1!80!black,
        right color=#1
      }
    },
    % سایه نرم با blur
    drop fuzzy shadow=black!30,
    left=8pt, right=8pt, top=8pt, bottom=8pt,
    after skip=8pt, before skip=8pt,
    % خط نازک تزئینی داخل
    underlay={\draw[#1!40, line width=0.3pt, dashed]
               (interior.south west) -- (interior.south east);},
  }
}

% ──── 🔵 کادر تعریف ──────────────────────────────────────────────────────────
\newtcolorbox{defbox}[1][]{
    bookbase=DefBoxFrame,
    colback=DefBoxBg,
    colframe=DefBoxFrame,
    coltitle=white,
    title={\rl{📖 #1}},
    boxed title style={colback=DefBoxFrame, rounded corners},
    % نوار جانبی
    borderline west={4pt}{-4pt}{DefBoxFrame!70},
    before upper={\parindent=1.5em},
}

% ──── 🟢 کادر نقل‌قول ────────────────────────────────────────────────────────
\newtcolorbox{quotebox}[1][]{
    bookbase=QuoteBoxFrame,
    colback=QuoteBoxBg,
    colframe=QuoteBoxFrame,
    coltitle=white,
    title={\rl{💬 #1}},
    boxed title style={colback=QuoteBoxFrame, rounded corners},
    borderline west={5pt}{-5pt}{QuoteBoxFrame},
    % علامت نقل‌قول تزئینی در پس‌زمینه
    watermark text={\textcolor{QuoteBoxFrame!15}{\fontsize{60}{60}\selectfont"}},
    before upper={\parindent=0em},
}

% ──── 🟡 کادر نکته تاریخی ───────────────────────────────────────────────────
\newtcolorbox{historybox}[1][]{
    bookbase=HistoryBoxFrame,
    colback=HistoryBoxBg,
    colframe=HistoryBoxFrame,
    coltitle=white,
    title={\rl{📜 #1}},
    boxed title style={colback=HistoryBoxFrame, rounded corners},
    % تزئین گوشه
    overlay={
      \node[anchor=south west, opacity=0.07,
            font=\fontsize{36}{36}\selectfont]
        at (interior.south west) {⏳};
    },
}

% ──── 🔴 کادر اختلافی ───────────────────────────────────────────────────────
\newtcolorbox{debatebox}[1][]{
    bookbase=DebateBoxFrame,
    colback=DebateBoxBg,
    colframe=DebateBoxFrame,
    coltitle=white,
    title={\rl{⚔️ #1}},
    boxed title style={colback=DebateBoxFrame, rounded corners},
    borderline={0.8pt}{-4pt}{DebateBoxFrame, dashed},
    overlay={
      \node[anchor=south east, opacity=0.06,
            font=\fontsize{40}{40}\selectfont]
        at (interior.south east) {⚔️};
    },
}

% ──── 🟣 کادر مقایسه ────────────────────────────────────────────────────────
\newtcolorbox{comparebox}[1][]{
    bookbase=CompareBoxFrame,
    colback=CompareBoxBg,
    colframe=CompareBoxFrame,
    coltitle=white,
    title={\rl{⚖️ #1}},
    boxed title style={colback=CompareBoxFrame, rounded corners},
    % دو نوار جانبی متقارن
    borderline west={3pt}{-3pt}{LeftRed!60},
    borderline east={3pt}{-3pt}{RightBlue!60},
}

% ──── 👤 کادر شخصیت ─────────────────────────────────────────────────────────
\newtcolorbox{personbox}[1][]{
    bookbase=PersonBoxFrame,
    colback=PersonBoxBg,
    colframe=PersonBoxFrame,
    coltitle=white,
    title={\rl{👤 #1}},
    boxed title style={colback=PersonBoxFrame, rounded corners},
    % هاشور سبک در پس‌زمینه
    interior style={
      top color=PersonBoxBg,
      bottom color=PersonBoxBg!85!PersonBoxFrame
    },
    before upper={\parindent=1.5em},
}

% ──── ⚠️ کادر هشدار ─────────────────────────────────────────────────────────
\newtcolorbox{warningbox}[1][]{
    bookbase=WarningBoxFrame,
    colback=WarningBoxBg,
    colframe=WarningBoxFrame,
    coltitle=white,
    title={\rl{⚠️ #1}},
    boxed title style={colback=WarningBoxFrame, rounded corners},
    borderline west={5pt}{-5pt}{WarningBoxFrame},
    overlay={
      \node[anchor=south east, opacity=0.06,
            font=\fontsize{40}{40}\selectfont]
        at (interior.south east) {⚠️};
    },
}

% ──── 📝 کادر خلاصه فصل ─────────────────────────────────────────────────────
\newtcolorbox{summarybox}{
    enhanced, breakable,
    colback=VeryLightGray,
    colframe=DarkGray,
    coltitle=white,
    fonttitle=\bfseries\large,
    title={\rl{📝 خلاصه فصل}},
    attach boxed title to top right={yshift=-2.5mm, xshift=-5mm},
    boxed title style={
      colback=DarkGray,
      rounded corners,
      % گرادیان روی تایتل
      interior style={left color=DarkGray!80, right color=DarkGray}
    },
    arc=4mm,
    boxrule=1.2pt,
    % کادر دوگانه
    frame style={
      top color=DarkGray!70,
      bottom color=DarkGray
    },
    drop fuzzy shadow=black!30,
    left=10pt, right=10pt, top=10pt, bottom=10pt,
    after skip=8pt, before skip=0pt,
    % خط تزئینی پایین
    underlay={
      \draw[GoldAccent, line width=1pt]
        (interior.south west) -- (interior.south east);
    },
}

% ──── ❓ کادر پرسش‌های تأملی ────────────────────────────────────────────────
\newtcolorbox{questionbox}{
    enhanced, breakable,
    colback=white,
    colframe=CenterGreen,
    coltitle=white,
    fonttitle=\bfseries\large,
    title={\rl{❓ پرسش‌های تأملی}},
    attach boxed title to top right={yshift=-2.5mm, xshift=-5mm},
    boxed title style={
      colback=CenterGreen,
      rounded corners,
      interior style={left color=CenterGreenDark, right color=CenterGreen}
    },
    arc=4mm,
    boxrule=1.2pt,
    borderline west={4pt}{-4pt}{CenterGreen},
    drop fuzzy shadow=black!20,
    left=10pt, right=10pt, top=10pt, bottom=10pt,
    after skip=0pt, before skip=8pt,
    underlay={
      \draw[CenterGreen!40, line width=0.4pt, dashed]
        ([xshift=4pt]interior.north west)
        -- ([xshift=4pt]interior.south west);
    },
}

% ══════════════════════════════════════════════════════════════════════════════
%                   🎨 استایل‌های TikZ
% ══════════════════════════════════════════════════════════════════════════════

\tikzstyle{timenode}=[rectangle, rounded corners=3pt,
    minimum width=2.5cm, minimum height=0.8cm,
    text centered, font=\small]
\tikzstyle{leftevent}=[timenode, fill=LeftRedLight,
    draw=LeftRed, line width=1pt, text=white]
\tikzstyle{rightevent}=[timenode, fill=RightBlueLight,
    draw=RightBlue, line width=1pt, text=white]
\tikzstyle{centerevent}=[timenode, fill=CenterGreenLight,
    draw=CenterGreen, line width=1pt, text=white]
\tikzstyle{neutralevent}=[timenode, fill=LightGray,
    draw=DarkGray, line width=1pt, text=DarkGray]

\tikzstyle{concept}=[ellipse, minimum width=2.5cm,
    minimum height=1.2cm, text centered, font=\small\bfseries]
\tikzstyle{leftconcept}=[concept, fill=LeftRedLight!50,
    draw=LeftRed, line width=1.5pt, text=LeftRedDark]
\tikzstyle{rightconcept}=[concept, fill=RightBlueLight!50,
    draw=RightBlue, line width=1.5pt, text=RightBlueDark]
\tikzstyle{centerconcept}=[concept, fill=CenterGreenLight!50,
    draw=CenterGreen, line width=1.5pt, text=CenterGreenDark]

\tikzstyle{thickarrow}=[{-Stealth[length=8pt,width=6pt]},line width=2pt]
\tikzstyle{thinarrow}=[{-Stealth[length=5pt,width=4pt]},line width=1pt]
\tikzstyle{dashedarrow}=[{-Stealth[length=5pt,width=4pt]},line width=1pt,dashed]

% ══════════════════════════════════════════════════════════════════════════════
%            📐 استایل فصل با TikZ (جایگزین titleformat ساده)
% ══════════════════════════════════════════════════════════════════════════════

% ─── Part ────────────────────────────────────────────────────────────────────
\titleformat{\part}[display]
    {\centering\Huge\bfseries}
    {%
      % نوار رنگی پشت شماره بخش
      \begin{tikzpicture}[remember picture]
        \fill[PartBg, rounded corners=4pt]
             (-0.5\linewidth, -0.3cm) rectangle (0.5\linewidth, 0.9cm);
        \draw[GoldAccent, line width=1.5pt, rounded corners=4pt]
             (-0.5\linewidth, -0.3cm) rectangle (0.5\linewidth, 0.9cm);
        \node[text=GoldAccent, font=\Large\bfseries]
             at (0, 0.3cm) {\rl{بخش \thepart}};
      \end{tikzpicture}%
    }
    {20pt}
    {\textcolor{DarkGray}}
\titlespacing*{\part}{0pt}{50pt}{40pt}

% ─── Chapter  ← BUG FIX: outer [] protected with braces ─────────────────────
\titleformat{\chapter}[display]
    {%
      \normalfont\huge\bfseries\raggedleft
      % پس‌زمینه نوار فصل (overlay)
      \begin{tikzpicture}[remember picture, overlay]
        % نوار بالای صفحه
        \fill[ChapterBandColor]
          ([yshift=-2pt]current page.north west)
          rectangle
          ([yshift=-8pt]current page.north east);
        % مستطیل شماره فصل
        \fill[ChapterBandColor, rounded corners=3pt]
          ([xshift=1.5cm, yshift=-3cm]current page.north east)
          rectangle
          ([xshift=4cm, yshift=-1.5cm]current page.north east);
        \node[text=ChapterNumColor,
              font=\fontsize{28}{28}\bfseries,
              anchor=center]
          at ([xshift=2.75cm, yshift=-2.25cm]current page.north east)
          {\thechapter};
      \end{tikzpicture}%
    }
    {\textcolor{MediumGray}{\Large\rl{فصل \thechapter}}}
    {10pt}
    {\textcolor{DarkGray}}
    % ↓ BUG FIX: braces around \titlerule[2pt] prevent stray ]
    [\vspace{4pt}%
     \ornamentalrule%
     \vspace{2pt}]
\titlespacing*{\chapter}{0pt}{-20pt}{30pt}

% ─── Section ─────────────────────────────────────────────────────────────────
\titleformat{\section}
    {\normalfont\Large\bfseries\color{RightBlue}}
    {%
      % شماره با پس‌زمینه کوچک
      \colorbox{RightBlue!12}{%
        \textcolor{RightBlue}{\thesection}%
      }%
    }
    {10pt}
    {}
    % خط کوتاه تزئینی زیر عنوان
    [\vspace{2pt}%
     {\color{RightBlue!40}\leaders\hrule height 0.5pt\hfill\kern0pt}%
     \vspace{1pt}]
\titlespacing*{\section}{0pt}{25pt}{15pt}

% ─── Subsection ──────────────────────────────────────────────────────────────
\titleformat{\subsection}
    {\normalfont\large\bfseries\color{DarkGray}}
    {\colorbox{DarkGray!10}{\textcolor{DarkGray}{\thesubsection}}}
    {8pt}
    {}
\titlespacing*{\subsection}{0pt}{20pt}{10pt}

% ─── Subsubsection ───────────────────────────────────────────────────────────
\titleformat{\subsubsection}
    {\normalfont\normalsize\bfseries\color{MediumGray}}
    {\thesubsubsection}
    {6pt}
    {}
\titlespacing*{\subsubsection}{0pt}{15pt}{8pt}

% ══════════════════════════════════════════════════════════════════════════════
%              📄 هدر و فوتر با TikZ (غنی‌شده)
% ══════════════════════════════════════════════════════════════════════════════

\pagestyle{fancy}
\fancyhf{}

% ─── هدر ─────────────────────────────────────────────────────────────────────
\fancyhead[RO]{\small\textcolor{DarkGray}{\leftmark}}
\fancyhead[LE]{\small\textcolor{DarkGray}{\rightmark}}

% شماره صفحه با دایره تزئینی
\fancyhead[LO]{%
  \begin{tikzpicture}[baseline=-0.6ex]
    \fill[ChapterBandColor!85] (0,0) circle (9pt);
    \node[text=white, font=\small\bfseries] at (0,0) {\thepage};
  \end{tikzpicture}%
}
\fancyhead[RE]{%
  \begin{tikzpicture}[baseline=-0.6ex]
    \fill[ChapterBandColor!85] (0,0) circle (9pt);
    \node[text=white, font=\small\bfseries] at (0,0) {\thepage};
  \end{tikzpicture}%
}

% خط هدر طلایی دو-لایه
\renewcommand{\headrulewidth}{0pt}   % غیرفعال‌سازی پیش‌فرض
\renewcommand{\headrule}{%
  \vspace{-4pt}%
  {\color{GoldAccent}\rule{\headwidth}{1.5pt}}\\[-2pt]
  {\color{ChapterBandColor!30}\rule{\headwidth}{0.4pt}}%
}

% فوتر
\fancyfoot[C]{}
\renewcommand{\footrulewidth}{0pt}

% ─── صفحات ساده (ابتدای فصل) ──────────────────────────────────────────────
\fancypagestyle{plain}{
    \fancyhf{}
    \fancyfoot[C]{%
      \begin{tikzpicture}[baseline=-0.6ex]
        \draw[GoldAccent, line width=0.8pt]
             (-1cm, 0) -- (-0.25cm, 0);
        \fill[GoldAccent] (0,0) circle (3pt);
        \node[text=DarkGray, font=\small\bfseries]
             at (0,0) {};
        \node[text=DarkGray, font=\small]
             at (0,-0.5) {\thepage};
        \draw[GoldAccent, line width=0.8pt]
             (0.25cm, 0) -- (1cm, 0);
      \end{tikzpicture}%
    }
    \renewcommand{\headrulewidth}{0pt}
}

% ══════════════════════════════════════════════════════════════════════════════
%                   📋 تنظیمات فهرست مطالب (غنی‌شده)
% ══════════════════════════════════════════════════════════════════════════════

\setcounter{tocdepth}{2}
\setcounter{secnumdepth}{3}

\titlecontents{part}[0pt]
    {\addvspace{20pt}\bfseries\large\color{ChapterBandColor}}
    {{\colorbox{ChapterBandColor!15}{\textcolor{ChapterBandColor}{بخش \thecontentslabel}}}\quad}
    {}
    {\hfill\contentspage}

\titlecontents{chapter}[15pt]
    {\addvspace{10pt}\bfseries\color{DarkGray}}
    {{\textcolor{GoldDark}{فصل \thecontentslabel:}}\quad}
    {}
    {\titlerule*[6pt]{\textcolor{GoldAccent!60}{·}}\contentspage}

\titlecontents{section}[35pt]
    {\addvspace{3pt}\color{MediumGray}}
    {\thecontentslabel\quad}
    {}
    {\titlerule*[8pt]{.}\contentspage}

\titlecontents{subsection}[55pt]
    {\color{MediumGray!80}}
    {\thecontentslabel\quad}
    {}
    {\titlerule*[8pt]{.}\contentspage}

% ══════════════════════════════════════════════════════════════════════════════
%                   🔧 دستورات سفارشی
% ══════════════════════════════════════════════════════════════════════════════

\newcommand{\fa}[1]{{\rl{#1}}}
\newcommand{\keyword}[1]{\textbf{\textcolor{RightBlue}{#1}}}
\newcommand{\leftkw}[1]{\textbf{\textcolor{LeftRed}{#1}}}
\newcommand{\rightkw}[1]{\textbf{\textcolor{RightBlue}{#1}}}
\newcommand{\centerkw}[1]{\textbf{\textcolor{CenterGreen}{#1}}}
\newcommand{\person}[1]{\textbf{#1}}
\newcommand{\book}[1]{\textit{«#1»}}
\newcommand{\term}[1]{\textbf{\textcolor{DarkGray}{#1}}}
\newcommand{\yearmark}[1]{\textbf{\textcolor{GoldDark}{#1}}}
\newcommand{\pmark}{\textcolor{LeftRed}{$\bullet$}}
\newcommand{\rmark}{\textcolor{RightBlue}{$\blacktriangleright$}}
\newcommand{\mediumspace}{\vspace{6pt}}
\newcommand{\largespace}{\vspace{12pt}}
\newcommand{\en}[1]{\lr{#1}}
\newcommand{\enfootnote}[1]{\footnote{\lr{#1}}}
\newcommand{\enref}[1]{\lr{\cite{#1}}}

% ──── جدیدها ────────────────────────────────────────────────────────────────
% خط جداکننده تزئینی (جایگزین \mediumspace در فاصله‌های مهم)
\newcommand{\bookrule}{%
  \vspace{4pt}\ornamentalrule\vspace{4pt}%
}
% تأکید سیاسی ترکیبی
\newcommand{\polmark}[1]{%
  \tikz[baseline=-0.5ex]{%
    \fill[GoldAccent!20, rounded corners=2pt]
         (0,-1pt) rectangle (\widthof{#1}+4pt, 9pt);
    \node[inner sep=2pt] at (0.5\widthof{#1}+2pt, 4pt) {#1};
  }%
}

% ══════════════════════════════════════════════════════════════════════════════
%                   📚 xepersian باید آخرین پکیج باشد
% ══════════════════════════════════════════════════════════════════════════════

\usepackage{xepersian}
\settextfont{Vazirmatn}
\setdigitfont{Vazirmatn}
%\PersianFootnotes